\documentclass[12pt,a4paper]{report}

\usepackage[utf8]{inputenc}
\usepackage[T5]{fontenc}
\usepackage[vietnamese]{babel}
\usepackage{lmodern}
\usepackage{textcomp}
\usepackage{upquote}
\usepackage{microtype}


\usepackage{amsmath}
\usepackage{amssymb}
\usepackage{mathtools}


\usepackage[a4paper,margin=2.4cm]{geometry}
\usepackage{titlesec}
\usepackage{fancyhdr}
\usepackage{enumitem}
\usepackage{graphicx}
\usepackage{xcolor}
\usepackage[most]{tcolorbox}
\usepackage{listings}
\usepackage{booktabs}
\usepackage{longtable}
\usepackage{array}
\usepackage{hyperref}
\usepackage{caption}
\usepackage{needspace}
\usepackage{float}
\definecolor{maintext}{RGB}{0,0,0} 

\definecolor{codebg}{RGB}{246,248,250}
\definecolor{codeframe}{RGB}{210,214,220}
\definecolor{kw}{RGB}{0,92,197}
\definecolor{str}{RGB}{163,21,21}
\definecolor{cmt}{RGB}{106,115,125}
\definecolor{num}{RGB}{0,110,90}
\definecolor{boxinfo}{RGB}{235,246,255}
\definecolor{boxinfoframe}{RGB}{84,150,220}
\definecolor{boxwarn}{RGB}{255,244,229}
\definecolor{boxwarnframe}{RGB}{214,140,40}
\definecolor{boxtip}{RGB}{234,250,240}
\definecolor{boxtipframe}{RGB}{50,160,90}
\definecolor{boxex}{RGB}{247,240,255}
\definecolor{boxexframe}{RGB}{130,80,200}
\definecolor{titleblue}{RGB}{20,60,110}


\titleformat{\chapter}[display]
  {\normalfont\Large\bfseries\color{titleblue}}
  {\filright\Large Chương \thechapter}
  {0.6em}{\Large\bfseries}
\titlespacing*{\chapter}{0pt}{20pt}{20pt}

\titleformat{\section}
  {\normalfont\Large\bfseries\color{titleblue}}
  {\thesection}{0.6em}{}
\titleformat{\subsection}
  {\normalfont\large\bfseries\color{titleblue!80!black}}
  {\thesubsection}{0.6em}{}
\titleformat{\subsubsection}
  {\normalfont\normalsize\bfseries\color{titleblue!70!black}}
  {\thesubsubsection}{0.6em}{}

\pagestyle{fancy}
\fancyhf{}
\renewcommand{\headrulewidth}{0.4pt}
\lhead{\small\textit{SymPy \& Cấp số nhân}}
\rhead{\small\thepage}


\lstdefinelanguage{PythonPlus}[]{Python}{
  morekeywords=[2]{sympy, symbols, Eq, solve, nonlinsolve, solveset, limit, oo,
                    simplify, subs, evalf, Rational, init_printing, latex,
                    lambdify, plot, expand, factor, nsimplify, sqrt, Abs, S,
                    SymPy, sin, cos, pi, E, ln, exp, log},
  keywordstyle=[2]\color{num}\bfseries,
}

\lstset{
  language=PythonPlus,
  backgroundcolor=\color{codebg},
  rulecolor=\color{codeframe},
  frame=single,
  framesep=6pt,
  framerule=0.6pt,
  basicstyle=\ttfamily\small,
  keywordstyle=\color{kw}\bfseries,
  stringstyle=\color{str},
  commentstyle=\color{cmt}\itshape,
  numbers=left,
  numberstyle=\tiny\color{cmt},
  numbersep=8pt,
  showstringspaces=false,
  breaklines=true,
  breakatwhitespace=false,
  postbreak=\mbox{\textcolor{red}{$\hookrightarrow$}\space}, % Ngắt dòng đẹp
  tabsize=4,
  columns=flexible,
  xleftmargin=1em,
  xrightmargin=0.3em,
  aboveskip=8pt,
  belowskip=8pt,
  captionpos=b,
}

\lstdefinestyle{outstyle}{
  backgroundcolor=\color{white},
  rulecolor=\color{codeframe},
  frame=single,
  framesep=6pt,
  basicstyle=\ttfamily\small\color{black!75},
  breaklines=true,
  numbers=none,
  xleftmargin=1em,
}


\newtcolorbox{infobox}[1][]{
  colback=boxinfo, colframe=boxinfoframe, boxrule=0.6pt, arc=2pt,
  left=7pt, right=7pt, top=5pt, bottom=5pt,
  title={Giải thích thêm}, fonttitle=\bfseries\small,
  coltitle=boxinfoframe!40!black, #1
}
\newtcolorbox{warnbox}[1][]{
  colback=boxwarn, colframe=boxwarnframe, boxrule=0.6pt, arc=2pt,
  left=7pt, right=7pt, top=5pt, bottom=5pt,
  title={Lưu ý / Lỗi thường gặp}, fonttitle=\bfseries\small,
  coltitle=boxwarnframe!30!black, #1
}
\newtcolorbox{tipbox}[1][]{
  colback=boxtip, colframe=boxtipframe, boxrule=0.6pt, arc=2pt,
  left=7pt, right=7pt, top=5pt, bottom=5pt,
  title={Mẹo thực hành}, fonttitle=\bfseries\small,
  coltitle=boxtipframe!30!black, #1
}
\newtcolorbox{exbox}[1][]{
  colback=boxex, colframe=boxexframe, boxrule=0.6pt, arc=2pt,
  left=7pt, right=7pt, top=5pt, bottom=5pt,
  title={Bài tập tự luyện}, fonttitle=\bfseries\small,
  coltitle=boxexframe!30!black, #1
}
\newtcolorbox{syntaxbox}[2][]{
  enhanced, colback=white, colframe=titleblue!60!black,
  boxrule=0.6pt, arc=2pt, left=7pt, right=7pt, top=5pt, bottom=5pt,
  title={Cú pháp: \texttt{#2}}, fonttitle=\bfseries\ttfamily\small,
  coltitle=white, colbacktitle=titleblue!70!black, #1
}

\newcommand{\pycmd}[1]{\texttt{\color{kw}#1}}
\newcommand{\file}[1]{\texttt{#1}}

\setlist[itemize]{leftmargin=1.4em, itemsep=2pt, topsep=3pt}
\setlist[enumerate]{leftmargin=1.6em, itemsep=3pt, topsep=3pt}

\hypersetup{
  colorlinks=true,
  linkcolor=titleblue,
  urlcolor=blue!60!black,
  citecolor=titleblue,
}


\title{\textbf{\color{titleblue}Hướng dẫn sử dụng SymPy \\ để giải các bài toán Cấp số nhân}\\[0.4em]}
\author{Sinh viên thực hiện}
\date{\today}


\begin{document}
\color{maintext} 

\begin{titlepage}
    \centering
    
    \vspace*{0.1cm}
    {\large \textbf{TRƯỜNG ĐẠI HỌC SƯ PHẠM -- ĐẠI HỌC HUẾ}} \\[0.2cm]
    {\large \textbf{KHOA TOÁN HỌC}}
    
    \vfill
    
    
        \includegraphics[width=4cm]{307101407_585748093337109_1610541625467017855_n.jpg}
    
    
    \vfill
    
    
    {\Large \textbf{HƯỚNG DẪN SỬ DỤNG SYMPY TRONG GIẢI QUYẾT CÁC VẤN ĐỀ CẤP SỐ NHÂN}}
    
    \vfill
    
    \begin{flushleft}
        \hspace{3.5cm} 
        \renewcommand{\arraystretch}{1.3} 
        \begin{tabular}{l @{\hspace{0.3cm}} c @{\hspace{0.3cm}} l}
            \textbf{Sinh viên thực hiện} & : & Trần Hoàng Vũ \\
            \textbf{Mã sinh viên} & : & 23S7010010 \\
            \textbf{Lớp} & : & 231012T \\
            \textbf{Khoá} & : & 2023 -- 2027 \\
            \textbf{Giảng viên hướng dẫn} & : & Nguyễn Đăng Minh Phúc \\
        \end{tabular}
    \end{flushleft}
    
    \vfill
    
    {\textbf{Thành phố Huế, tháng 9 năm 2026}}
    \vspace*{0.5cm}
    
\end{titlepage}



\newpage
\tableofcontents
\newpage


\chapter*{Mở đầu: Đặt vấn đề và cơ sở lý thuyết}
\addcontentsline{toc}{chapter}{Mở đầu: Đặt vấn đề và cơ sở lý thuyết}

\section*{0.1. Nhắc lại lý thuyết Cấp số nhân}
\addcontentsline{toc}{section}{0.1. Nhắc lại lý thuyết Cấp số nhân}

Một dãy số $(u_n)$ được gọi là \textbf{cấp số nhân} nếu kể từ số hạng
thứ hai, mỗi số hạng đều bằng số hạng đứng ngay trước nó nhân với một
số không đổi $q$ (gọi là \textbf{công bội}):
\[
u_{n+1} = u_n \cdot q, \qquad \forall n \geq 1.
\]

Các đại lượng cơ bản cần nhớ:

\begin{itemize}
  \item $u_1$: số hạng đầu tiên của dãy.
  \item $q$: công bội (tỉ số giữa hai số hạng liên tiếp, $q \neq 0$).
  \item \textbf{Số hạng tổng quát:}
  \[
  u_n = u_1 \cdot q^{\,n-1}.
  \]
  \item \textbf{Tính chất ba số hạng liên tiếp:}
  $u_k^2 = u_{k-1} \cdot u_{k+1}$ (với $k \geq 2$).
  \item \textbf{Tổng $n$ số hạng đầu tiên} $S_n = u_1 + u_2 + \dots + u_n$:
  \[
  S_n =
  \begin{cases}
    n \cdot u_1, & \text{nếu } q = 1,\\[4pt]
    u_1 \cdot \dfrac{1 - q^n}{1 - q}, & \text{nếu } q \neq 1.
  \end{cases}
  \]
  \item \textbf{Cấp số nhân lùi vô hạn:} khi $|q| < 1$ và $n \to \infty$,
  tổng $S_n$ tiến tới một giá trị hữu hạn:
  \[
  S = \lim_{n \to \infty} S_n = \frac{u_1}{1 - q}.
  \]
\end{itemize}

\begin{infobox}
Bốn ``khối kiến thức'' toàn bộ tài liệu sẽ xoay quanh: \emph{số hạng đầu
$u_1$}, \emph{công bội $q$}, \emph{số hạng tổng quát $u_n$}, và
\emph{tổng $S_n$ (hoặc $S$ khi lùi vô hạn)}. SymPy sẽ giúp ta thao tác
với các công thức này như những \textbf{biểu thức đại số biết suy luận},
thay vì chỉ là những con số cụ thể.
\end{infobox}

\section*{0.2. SymPy là gì? Vì sao dùng SymPy để giải toán?}
\addcontentsline{toc}{section}{0.2. SymPy là gì? Vì sao dùng SymPy?}

\textbf{SymPy} là một thư viện \textbf{Toán học ký hiệu} (Symbolic
Mathematics) viết bằng Python, thuộc nhóm phần mềm gọi là \textbf{CAS}
(\emph{Computer Algebra System} -- Hệ thống đại số máy tính).

Khác biệt cốt lõi so với tính toán số học thông thường:

\begin{center}
\begin{tabular}{@{}p{0.46\textwidth}p{0.46\textwidth}@{}}
\toprule
\textbf{Số học (numeric)} & \textbf{Ký hiệu (SymPy)}\\
\midrule
Biến phải gán một giá trị số cụ thể, ví dụ \texttt{q = 0.5}. &
Biến có thể là một \textbf{ký hiệu trừu tượng} $q$, chưa cần biết giá trị.\\[4pt]
Kết quả là số gần đúng, có thể mất chính xác (sai số làm tròn). &
Kết quả là biểu thức \textbf{chính xác tuyệt đối} (phân số, căn thức, $\pi$, \dots).\\[4pt]
Muốn đổi bài toán phải viết lại công thức tính bằng tay. &
Có thể \textbf{giải phương trình, rút gọn, lấy giới hạn} ngay trên biểu thức.\\[4pt]
Phù hợp tính số liệu lớn, xử lý ảnh, mô phỏng. &
Phù hợp giải toán đại số, chứng minh công thức, kiểm tra đáp án.\\
\bottomrule
\end{tabular}
\end{center}

Nói cách khác, với SymPy ta có thể gõ đúng công thức toán học
$u_n = u_1 q^{n-1}$ vào Python gần như y hệt cách viết trên giấy, rồi
nhờ máy tính thay số, giải phương trình, hoặc rút gọn giúp mình -- giảm
tối đa sai sót tính tay.

\begin{tipbox}
Tài liệu dùng \textbf{Jupyter Notebook} làm môi trường thực hành (theo
đúng định dạng đã thống nhất trong kế hoạch dự án), vì Notebook cho
phép xen kẽ giữa \emph{lý thuyết trình bày bằng Markdown/LaTeX} và
\emph{mã Python chạy thử ngay lập tức}. Tuy nhiên mọi đoạn mã trong tài
liệu đều chạy được y hệt nếu bạn dùng file \texttt{.py} thông thường.
\end{tipbox}

\section*{0.3. Lộ trình tài liệu}
\addcontentsline{toc}{section}{0.3. Lộ trình tài liệu}

Tài liệu gồm 9 chương, độ khó tăng dần:

\begin{center}
\begin{tabular}{@{}cl@{}}
\toprule
\textbf{Chương} & \textbf{Nội dung}\\
\midrule
1 & Cài đặt môi trường, khai báo biến ký hiệu, in đẹp\\
2 & Các lệnh SymPy cơ bản: \texttt{subs}, \texttt{evalf}, \texttt{simplify}\\
3 & Dịch công thức CSN sang SymPy (\texttt{Eq}, \texttt{solve})\\
4 & Dạng 1: Tính số hạng và tổng khi biết $u_1, q$\\
5 & Dạng 2: Giải hệ phương trình tìm $u_1, q$, lọc nghiệm\\
6 & Dạng 3: Cấp số nhân lùi vô hạn với \texttt{limit()}\\
7 & Ứng dụng thực tế: lãi kép, vi khuẩn, dân số\\
8 & Trực quan hóa bằng \texttt{matplotlib} + \texttt{lambdify}\\
9 & Đóng gói module \texttt{solver.py} và kiểm thử\\
\bottomrule
\end{tabular}
\end{center}

\newpage

% =====================================================================
\chapter{Cài đặt môi trường và khai báo biến ký hiệu}

\section{Cài đặt môi trường làm việc}

\textbf{Yêu cầu hệ thống:} Python phiên bản 3.9 trở lên.

Mở cửa sổ dòng lệnh (Terminal / Command Prompt) và cài ba thư viện cần
thiết:

\begin{lstlisting}[language=bash,caption={Cài đặt thư viện qua pip},numbers=none]
pip install sympy jupyter matplotlib
\end{lstlisting}

\begin{itemize}
  \item \pycmd{sympy}: thư viện toán học ký hiệu -- trung tâm tài liệu.
  \item \pycmd{jupyter}: môi trường Notebook tương tác.
  \item \pycmd{matplotlib}: thư viện vẽ đồ thị (dùng ở Chương 8).
\end{itemize}

Khởi động Jupyter Notebook:

\begin{lstlisting}[language=bash,numbers=none]
jupyter notebook
\end{lstlisting}

\begin{tipbox}
Kiểm tra SymPy đã cài đúng chưa bằng đoạn code sau. Nếu in ra số phiên
bản (không báo lỗi) là thành công:

\begin{lstlisting}
import sympy as sp
print(sp.__version__)   # e.g. 1.12
\end{lstlisting}
\end{tipbox}

\section{Nạp thư viện và quy ước đặt tên}

\begin{lstlisting}
import sympy as sp
\end{lstlisting}

\begin{syntaxbox}{import tên\_thư\_viện as bí\_danh}
\pycmd{import sympy as sp} nghĩa là: nạp toàn bộ thư viện
\texttt{sympy}, từ giờ gọi tắt bằng tên \texttt{sp}. Nhờ vậy thay vì gõ
\texttt{sympy.symbols(...)} dài dòng, ta chỉ cần
\texttt{sp.symbols(...)}. Đây là quy ước phổ biến trong cộng đồng
SymPy, tương tự \texttt{import numpy as np} hay
\texttt{import pandas as pd}.
\end{syntaxbox}

\begin{warnbox}
Nếu Python báo \texttt{ModuleNotFoundError: No module named 'sympy'},
nghĩa là bước cài đặt chưa thành công hoặc bạn đang chạy nhầm môi
trường Python. Hãy chạy lại \texttt{pip install sympy} trong đúng
terminal mà Jupyter đang sử dụng.
\end{warnbox}

Quy ước tên biến Python (bám sát sách giáo khoa Việt Nam):

\begin{center}
\begin{tabular}{@{}llll@{}}
\toprule
\textbf{Toán học} & \textbf{Biến Python} & \textbf{Ý nghĩa} & \textbf{Ví dụ}\\
\midrule
$u_1$ & \texttt{u1} & số hạng đầu & \texttt{u1 = 3}\\
$q$ & \texttt{q} & công bội & \texttt{q = 2}\\
$n$ & \texttt{n} & chỉ số/số thứ tự & \texttt{n = 5}\\
$u_n$ & \texttt{un} & số hạng tổng quát & \texttt{un = u1*q**(n-1)}\\
$S_n$ & \texttt{Sn} & tổng $n$ số hạng đầu & \texttt{Sn = ...}\\
\bottomrule
\end{tabular}
\end{center}

\section{Khai báo biến ký hiệu với \texttt{symbols()}}

\begin{syntaxbox}{sp.symbols('tên1 tên2 ...')}
Hàm \pycmd{symbols()} tạo ra các \textbf{đối tượng ký hiệu toán học}
(\texttt{Symbol}) -- những cái tên đại diện cho một đại lượng \emph{chưa
biết giá trị}, giống hệt cách ta viết $x$, $q$, $n$ trên giấy mà không
cần biết chúng bằng bao nhiêu.
\end{syntaxbox}

\begin{lstlisting}[caption={Khai báo các biến của bài toán cấp số nhân}]
import sympy as sp

# Declare symbolic variables used for geometric progression
u1, q, n = sp.symbols('u1 q n')

print(u1, q, n)
print(type(u1))
\end{lstlisting}

\begin{lstlisting}[style=outstyle,caption={Kết quả in ra}]
u1 q n
<class 'sympy.core.symbol.Symbol'>
\end{lstlisting}

\begin{infobox}
Chú ý sự khác nhau giữa \emph{tên biến Python} (\texttt{u1}) và
\emph{chuỗi ký tự} trong \texttt{symbols('u1 q n')}. Dòng lệnh
\texttt{u1, q, n = sp.symbols('u1 q n')} làm \textbf{hai việc cùng lúc}:
(1) yêu cầu SymPy tạo ba ký hiệu có tên hiển thị \texttt{"u1"},
\texttt{"q"}, \texttt{"n"}; (2) gán chúng cho ba biến Python cùng tên để
thao tác trong code. Tên biến Python và tên hiển thị không bắt buộc
phải trùng nhau (ví dụ \texttt{a = sp.symbols('u1')} vẫn hợp lệ), nhưng
đặt trùng tên giúp code dễ đọc hơn.
\end{infobox}

Một số biến thể hữu ích:

\begin{lstlisting}[caption={Các cách khai báo biến khác nhau}]
# Style 1: multiple symbols in one string, space-separated
a, b, c = sp.symbols('a b c')

# Style 2: comma-separated string -- same result
a, b, c = sp.symbols('a, b, c')

# Style 3: generate indexed symbols u1, u2, ..., u5
u = sp.symbols('u1:6')   # produces (u1, u2, u3, u4, u5)
print(u)

# Style 4: constrain domain -- useful for progressions
# q real, q nonzero (geometric ratio must not be zero)
q = sp.symbols('q', real=True, nonzero=True)
n = sp.symbols('n', positive=True, integer=True)
\end{lstlisting}

\begin{lstlisting}[style=outstyle,numbers=none]
(u1, u2, u3, u4, u5)
\end{lstlisting}

\begin{warnbox}
Khai báo thêm \texttt{real=True} (số thực), \texttt{positive=True}
(dương), \texttt{integer=True} (số nguyên), \texttt{nonzero=True}
(khác 0) \textbf{không bắt buộc}, nhưng \textbf{rất nên làm} với bài
toán cấp số nhân: vì $q$ phải khác 0 theo định nghĩa, và $n$ luôn là
số nguyên dương. Nếu bỏ qua, khi giải phương trình bậc cao SymPy có
thể trả về thêm nghiệm phức hoặc nghiệm âm không phù hợp ngữ cảnh.
\end{warnbox}

\section{Định dạng hiển thị đẹp với \texttt{init\_printing()}}

\begin{syntaxbox}{sp.init\_printing()}
Bật chế độ hiển thị công thức toán học dưới dạng \textbf{đẹp}
(Unicode hoặc LaTeX-style). Chỉ cần gọi \textbf{một lần} ở đầu
Notebook, mọi kết quả in ra sau đó sẽ tự động được định dạng.
\end{syntaxbox}

\begin{lstlisting}[caption={Bật chế độ in đẹp}]
import sympy as sp
sp.init_printing()   # enable pretty-print for entire notebook

u1, q, n = sp.symbols('u1 q n', real=True)
bieu_thuc = u1 * q**(n - 1)
bieu_thuc   # in Jupyter: shown as mathematical formula
\end{lstlisting}

Khi chạy trong Jupyter, ô kết quả hiển thị:
\[
u_1 \cdot q^{\,n-1}
\]
thay vì chuỗi thô \texttt{u1*q**(n - 1)}.

\begin{infobox}
Nếu không dùng Jupyter (chạy file \texttt{.py} thường qua terminal),
\texttt{init\_printing()} sẽ tự chuyển sang hiển thị bằng ký tự Unicode
thay vì công thức đồ họa, nhưng vẫn dễ đọc hơn nhiều so với mặc định.
\end{infobox}

\section{Thiết lập công thức CSN dưới dạng SymPy}

\begin{lstlisting}[caption={Biểu thức số hạng tổng quát và tổng $S_n$}]
import sympy as sp
sp.init_printing()

u1, q, n = sp.symbols('u1 q n', real=True)

# General term: u_n = u1 * q^(n-1)
un_formula = u1 * q**(n - 1)

# Sum of n terms (case q != 1):
# S_n = u1 * (1 - q**n) / (1 - q)
Sn_formula = u1 * (1 - q**n) / (1 - q)

print("Formula u_n:", un_formula)
print("Formula S_n:", Sn_formula)
\end{lstlisting}

\begin{lstlisting}[style=outstyle,caption={Kết quả in ra}]
Formula u_n: u1*q**(n - 1)
Formula S_n: u1*(1 - q**n)/(1 - q)
\end{lstlisting}

\begin{tipbox}
Từ đây trở đi, \texttt{un\_formula} và \texttt{Sn\_formula} không phải
là những con số cố định, mà là ``khuôn công thức'' sống động: bạn có
thể thay số vào (Chương 4), gộp với một phương trình khác để giải hệ
(Chương 5), hoặc lấy giới hạn khi $n \to \infty$ (Chương 6) -- tất cả
chỉ với vài dòng lệnh.
\end{tipbox}

\begin{exbox}
\textbf{Bài tập 1.1.} Khai báo các biến ký hiệu $a, r, m$ và tạo biểu
thức $T_m = a \cdot r^m$.\\[2pt]
\textbf{Bài tập 1.2.} Khai báo biến $k$ là số nguyên dương và $p$ là số
thực dương. In ra biểu thức $p^k$.\\[2pt]
\textbf{Bài tập 1.3.} Dùng \texttt{symbols} tạo một lượt các biến
$v_1, v_2, \ldots, v_8$ (gợi ý: cú pháp \texttt{'v1:9'}).
\end{exbox}

\newpage

% =====================================================================
\chapter{Các lệnh SymPy cơ bản}

Chương này giới thiệu từng lệnh SymPy riêng lẻ -- là ``viên gạch'' xây
nên mọi lời giải cho bài toán cấp số nhân ở các chương sau.

\section{Thay giá trị -- \texttt{.subs()}}

\begin{syntaxbox}{bieu\_thuc.subs(bien, gia\_tri)}
\texttt{.subs()} (viết tắt \emph{substitute} -- thay thế) nhận vào một
hoặc nhiều cặp ``biến -- giá trị'', trả về một biểu thức \textbf{mới}
với biến đã được thay bằng giá trị tương ứng. Biểu thức gốc \textbf{không
bị thay đổi} -- \texttt{.subs()} luôn trả về kết quả mới, cần gán vào
biến khác để lưu lại.
\end{syntaxbox}

\begin{lstlisting}[caption={Thay số cho nhiều biến cùng lúc}]
import sympy as sp
sp.init_printing()

u1, q, n = sp.symbols('u1 q n', real=True)
un_formula = u1 * q**(n - 1)

# Example: u1 = 3, q = 2, find u5
ket_qua = un_formula.subs({u1: 3, q: 2, n: 5})
print(ket_qua)
\end{lstlisting}

\begin{lstlisting}[style=outstyle,numbers=none]
48
\end{lstlisting}

\begin{infobox}
Cú pháp \texttt{.subs(\{u1: 3, q: 2, n: 5\})} dùng một \textbf{dictionary}
(từ điển) của Python để thay \emph{nhiều biến cùng lúc} -- mỗi cặp
\texttt{bien: gia\_tri} cách nhau bằng dấu phẩy. Bạn cũng có thể thay
từng biến một bằng \texttt{.subs(u1, 3).subs(q, 2).subs(n, 5)}, kết quả
hoàn toàn giống nhau. Cách dictionary gọn và rõ ràng hơn khi có nhiều
biến.
\end{infobox}

\section{Chuyển sang số thực -- \texttt{.evalf()}}

\begin{syntaxbox}{bieu\_thuc.evalf(so\_chu\_so)}
\texttt{.evalf()} (evaluate as float) ép một biểu thức đã có đầy đủ giá
trị số thành \textbf{số thập phân gần đúng}. Tham số (mặc định 15) quy
định số chữ số có nghĩa muốn hiển thị.
\end{syntaxbox}

\begin{lstlisting}[caption={Chuyển sang số thập phân}]
bieu_thuc = sp.sqrt(2) + sp.Rational(1, 3)
print(bieu_thuc)             # sqrt(2) + 1/3
print(bieu_thuc.evalf())     # 1.74754689571453
print(bieu_thuc.evalf(20))   # 1.7475468957145320588
\end{lstlisting}

\begin{infobox}
\pycmd{sp.Rational(1, 3)} tạo phân số chính xác $\tfrac13$ thay vì số
thập phân xấp xỉ \texttt{0.333...}. Đây là thói quen \textbf{rất quan
trọng} khi làm việc với SymPy: gõ \texttt{1/3} bằng cú pháp Python
thường sẽ cho số thực gần đúng \texttt{0.3333333333333333} ngay từ đầu
(mất chính xác); còn \texttt{sp.Rational(1, 3)} giữ nguyên phân số cho
đến bước cuối khi bạn gọi \texttt{.evalf()}.
\end{infobox}

\section{Rút gọn, khai triển, phân tích nhân tử}

\begin{lstlisting}[caption={Các lệnh biến đổi biểu thức}]
q = sp.symbols('q', real=True)

expr = (q**3 - 1) / (q - 1)
print("simplify:", sp.simplify(expr))     # q**2 + q + 1
print("expand:  ", sp.expand((q + 1)**3)) # q**3 + 3*q**2 + 3*q + 1
print("factor:  ", sp.factor(q**2 - 1))   # (q - 1)*(q + 1)
\end{lstlisting}

\begin{itemize}
  \item \pycmd{sp.simplify(expr)}: cố gắng tối giản biểu thức (dùng
  nhiều thuật toán nội bộ).
  \item \pycmd{sp.expand(expr)}: khai triển ngoặc.
  \item \pycmd{sp.factor(expr)}: phân tích thành nhân tử.
\end{itemize}

\section{Phát biểu phương trình -- \texttt{Eq()}}

\begin{warnbox}
Đây là lỗi phổ biến nhất với người mới học SymPy: trong Python, dấu
\texttt{==} chỉ dùng để \textbf{so sánh} (trả về \texttt{True}/\texttt{False}
ngay lập tức), \textbf{không} dùng để biểu diễn một phương trình toán
học chưa giải. Muốn viết phương trình như $u_2 + u_4 = 30$ để SymPy xử
lý, bắt buộc phải dùng hàm \texttt{sp.Eq(ve\_trai, ve\_phai)}.
\end{warnbox}

\begin{syntaxbox}{sp.Eq(ve\_trai, ve\_phai)}
Tạo ra một \textbf{đối tượng phương trình} (đẳng thức biểu tượng). Đối
tượng này chưa được ``giải'' -- nó chỉ ghi lại mối quan hệ, tương tự
việc bạn viết $u_2 + u_4 = 30$ ra giấy trước khi bắt đầu giải.
\end{syntaxbox}

\begin{lstlisting}[caption={Tạo và quan sát phương trình}]
u1, q = sp.symbols('u1 q', real=True)

pt = sp.Eq(u1 + u1*q, 6)
print(pt)              # Eq(q*u1 + u1, 6)
print(pt.lhs)          # q*u1 + u1  (ve trai)
print(pt.rhs)          # 6          (ve phai)
\end{lstlisting}

\section{Giải phương trình -- \texttt{solve()}}

\begin{syntaxbox}{sp.solve(pt, bien)}
Giải phương trình \texttt{pt} theo ẩn \texttt{bien}, trả về danh sách
các nghiệm. Với hệ nhiều phương trình, truyền vào danh sách phương
trình và danh sách ẩn:
\texttt{sp.solve([pt1, pt2], [x, y])}.
\end{syntaxbox}

\begin{lstlisting}[caption={Giải một phương trình}]
u1, q = sp.symbols('u1 q', real=True)
pt = sp.Eq(u1 + u1*q, 6)
print(sp.solve(pt, u1))     # [6/(q + 1)]
\end{lstlisting}

\begin{exbox}
\textbf{Bài tập 2.1.} Dùng \texttt{.subs()} tính $3 \cdot 2^9$ và
$5 \cdot 3^{4}$.\\[2pt]
\textbf{Bài tập 2.2.} Tính \texttt{sqrt(2) * sqrt(8)} -- kết quả ký hiệu
là gì? Chuyển sang số thập phân bằng \texttt{.evalf()}.\\[2pt]
\textbf{Bài tập 2.3.} Rút gọn phân thức
$\dfrac{q^4 - 1}{q - 1}$.\\[2pt]
\textbf{Bài tập 2.4.} Giải phương trình $q^2 - 5q + 6 = 0$.
\end{exbox}

\newpage

% =====================================================================
\chapter{Giải các dạng toán Cấp số nhân cơ bản}

\section{Dạng 1 -- Tìm số hạng thứ $n$ khi biết $u_1, q$}

Đây là dạng toán đơn giản nhất: đề bài đã cho sẵn $u_1$ và $q$, chỉ cần
\textbf{thay số} vào công thức bằng \texttt{.subs()}.

\begin{lstlisting}[caption={Ví dụ 1: Tìm $u_5$ khi $u_1 = 3, q = 2$}]
import sympy as sp
sp.init_printing()

u1, q, n = sp.symbols('u1 q n', real=True)
un_formula = u1 * q**(n - 1)

# u1 = 3, q = 2, find u5
ket_qua = un_formula.subs({u1: 3, q: 2, n: 5})
print(ket_qua)
\end{lstlisting}

\begin{lstlisting}[style=outstyle,numbers=none]
48
\end{lstlisting}

\noindent\textbf{Kiểm tra bằng tay:}
$u_5 = 3 \cdot 2^{4} = 3 \cdot 16 = 48$. Khớp.

\begin{lstlisting}[caption={Ví dụ 2: Tìm $u_{10}$ khi $u_1 = 1, q = 3$}]
ket_qua = un_formula.subs({u1: 1, q: 3, n: 10})
print(ket_qua)      # 19683
\end{lstlisting}

\section{Dạng 2 -- Tính tổng $S_n$ khi biết $u_1, q$}

\begin{lstlisting}[caption={Ví dụ 3: Tính $S_6$ khi $u_1 = 3, q = 2$}]
Sn_formula = u1 * (1 - q**n) / (1 - q)

S6 = Sn_formula.subs({u1: 3, q: 2, n: 6})
print(S6)
\end{lstlisting}

\begin{lstlisting}[style=outstyle,numbers=none]
189
\end{lstlisting}

\noindent\textbf{Kiểm tra bằng tay:}
$S_6 = 3 \cdot \dfrac{2^6 - 1}{2 - 1} = 3 \cdot 63 = 189$.

\begin{warnbox}
Nếu $q$ là một số \textbf{không nguyên} hoặc một \textbf{phân số}, kết
quả \texttt{.subs()} sẽ trả về dạng \textbf{phân số chính xác} (ví dụ
$\frac{189}{4}$) thay vì số thập phân gần đúng -- đây chính là ưu điểm
``tính toán ký hiệu''. Muốn xem giá trị thập phân, dùng thêm
\texttt{.evalf()}.
\end{warnbox}

\section{Ví dụ tổng hợp -- kết quả dạng phân số}

\begin{lstlisting}[caption={Tính $S_{10}$ khi $u_1 = 1, q = 1/2$}]
u1, q, n = sp.symbols('u1 q n', real=True)
Sn_formula = u1 * (1 - q**n) / (1 - q)

S10 = Sn_formula.subs({u1: 1, q: sp.Rational(1, 2), n: 10})
print("Exact fraction:", S10)       # 1023/512
print("Decimal (10 digits):", S10.evalf(10))  # 1.998046875
\end{lstlisting}

\section{Trường hợp đặc biệt $q = 1$}

\begin{warnbox}
Công thức $S_n = u_1 \dfrac{1-q^n}{1-q}$ chứa phép chia cho $(1 - q)$.
Khi $q = 1$, mẫu số bằng 0 -- thay trực tiếp $q=1$ vào công thức này
(kể cả trong SymPy) sẽ báo lỗi \texttt{ZeroDivisionError} hoặc trả về
\texttt{zoo} (ký hiệu ``vô cùng phức'' của SymPy), \textbf{không phải}
kết quả toán học đúng. Trường hợp $q = 1$ phải xử lý riêng bằng công
thức $S_n = n \cdot u_1$.
\end{warnbox}

\begin{lstlisting}[caption={Hàm tính Sn có xử lý riêng trường hợp q = 1}]
def tinh_Sn(u1_val, q_val, n_val):
    """Compute S_n, handling the special case q = 1."""
    if q_val == 1:
        return n_val * u1_val
    else:
        u1s, qs, ns = sp.symbols('u1 q n')
        cong_thuc = u1s * (1 - qs**ns) / (1 - qs)
        return cong_thuc.subs({u1s: u1_val, qs: q_val, ns: n_val})

print("q = 1 case:", tinh_Sn(5, 1, 10))   # 5*10 = 50
print("q = 3 case:", tinh_Sn(5, 3, 4))    # 200
\end{lstlisting}

\begin{lstlisting}[style=outstyle,caption={Kết quả}]
q = 1 case: 50
q = 3 case: 200
\end{lstlisting}

\begin{tipbox}
\textbf{Quy trình chuẩn cho Dạng 1:} (1) viết công thức tổng quát một
lần bằng ký hiệu; (2) dùng \texttt{.subs()} thay số theo đề bài; (3)
nếu cần số thập phân để báo cáo, gọi thêm \texttt{.evalf()}. Cách làm
này giúp tái sử dụng đúng một công thức cho hàng chục bài tập khác
nhau, chỉ đổi số liệu đầu vào.
\end{tipbox}

\begin{exbox}
\textbf{Bài tập 3.1.} Cho CSN có $u_1 = 2$, $q = 5$. Tìm $u_7$ và
$S_7$.\\[2pt]
\textbf{Bài tập 3.2.} Cho $u_1 = 1$, $q = \tfrac{1}{3}$. Tính $S_{20}$
dưới dạng phân số chính xác.\\[2pt]
\textbf{Bài tập 3.3.} Cho $u_1 = 100$, $q = 0.98$. Tính $u_{12}$ và
$S_{12}$ dưới dạng số thập phân (làm tròn 4 chữ số).
\end{exbox}

\newpage

% =====================================================================
\chapter{Dạng toán nâng cao -- Giải hệ phương trình tìm $u_1, q$}

\section{Bài toán đặt vấn đề}

Dạng toán này khó hơn: đề bài \textbf{không cho trực tiếp} $u_1$ và
$q$, mà cho các \textbf{mối liên hệ} giữa các số hạng, ví dụ:
\[
\begin{cases}
u_2 + u_4 = 30\\
u_3 + u_5 = 90
\end{cases}
\]
Yêu cầu: tìm $u_1$ và $q$. Ta cần hai công cụ: \texttt{Eq()} (đã học ở
Chương 2) và \texttt{solve()} / \texttt{nonlinsolve()}.

\section{Thiết lập hệ phương trình}

\begin{lstlisting}[caption={Thiết lập hệ phương trình từ đề bài}]
import sympy as sp
sp.init_printing()

u1, q = sp.symbols('u1 q', real=True)

# Build general terms for the required indices
u2 = u1 * q**(2 - 1)   # = u1*q
u3 = u1 * q**(3 - 1)   # = u1*q**2
u4 = u1 * q**(4 - 1)   # = u1*q**3
u5 = u1 * q**(5 - 1)   # = u1*q**4

# Given: u2 + u4 = 30  and  u3 + u5 = 90
pt1 = sp.Eq(u2 + u4, 30)
pt2 = sp.Eq(u3 + u5, 90)

print(pt1)
print(pt2)
\end{lstlisting}

\begin{lstlisting}[style=outstyle,caption={Kết quả}]
Eq(u1*q + u1*q**3, 30)
Eq(u1*q**2 + u1*q**4, 90)
\end{lstlisting}

\section{Giải hệ với \texttt{solve()}}

\begin{syntaxbox}{sp.solve([pt1, pt2], [bien1, bien2])}
Nhận vào \textbf{danh sách phương trình} (dấu ngoặc vuông) và
\textbf{danh sách ẩn số} cần tìm, trả về danh sách các \textbf{nghiệm}
thỏa mãn đồng thời tất cả phương trình.
\end{syntaxbox}

\begin{lstlisting}[caption={Giải hệ phương trình tìm u1, q}]
nghiem = sp.solve([pt1, pt2], [u1, q])
print(nghiem)
\end{lstlisting}

\begin{lstlisting}[style=outstyle,numbers=none]
[(10, 3), (10/9, -3)]
\end{lstlisting}

\begin{infobox}
Kết quả trả về là \textbf{danh sách các bộ nghiệm} (mỗi bộ trong ngoặc
đơn, theo đúng thứ tự \texttt{[u1, q]} đã khai báo). Hệ có \textbf{hai
bộ nghiệm}: $(u_1, q) = (10, 3)$ và $(u_1, q) = (\tfrac{10}{9}, -3)$.
Việc phương trình bậc cao cho \emph{nhiều nghiệm} là bình thường trong
đại số -- nhiệm vụ của người giải toán là dựa vào \textbf{điều kiện
thực tế của đề bài} để chọn ra nghiệm phù hợp.
\end{infobox}

\section{Giải hệ với \texttt{nonlinsolve()}}

\begin{syntaxbox}{sp.nonlinsolve([pt1, pt2], [bien1, bien2])}
Tương tự \texttt{solve()}, nhưng chuyên biệt cho \textbf{hệ phương
trình phi tuyến} (nonlinear system) -- có chứa ẩn với số mũ cao hơn 1,
hoặc tích các ẩn. Trả về kết quả dưới dạng \texttt{FiniteSet}.
\end{syntaxbox}

\begin{lstlisting}[caption={Giải hệ bằng nonlinsolve}]
nghiem_nonlin = sp.nonlinsolve([pt1, pt2], [u1, q])
print(nghiem_nonlin)
\end{lstlisting}

\begin{lstlisting}[style=outstyle,numbers=none]
{(10, 3), (10/9, -3)}
\end{lstlisting}

\begin{tipbox}
\textbf{Nên dùng hàm nào?} Với hệ CSN thông thường (2 ẩn $u_1, q$), cả
\texttt{solve()} và \texttt{nonlinsolve()} thường cho cùng kết quả --
\texttt{solve()} có cú pháp quen thuộc, kết quả dễ duyệt (\texttt{list}
các \texttt{tuple}). \texttt{nonlinsolve()} đáng tin cậy hơn khi hệ
\textbf{thực sự phức tạp}. Thực hành: thử \texttt{solve()} trước, nếu
báo không tìm được nghiệm mới chuyển sang \texttt{nonlinsolve()}.
\end{tipbox}

\section{Lọc nghiệm hợp lệ -- loại nghiệm phức}

Khi phương trình tìm $q$ có bậc cao, SymPy có thể trả về cả nghiệm phức
(chứa đơn vị ảo \texttt{sp.I}, ký hiệu $i$ với $i^2 = -1$).

\begin{lstlisting}[caption={Giải phương trình bậc ba và lọc nghiệm thực}]
q = sp.symbols('q')

# Solve q**3 = 8
pt = sp.Eq(q**3, 8)
tat_ca_nghiem = sp.solve(pt, q)
print("All roots:", tat_ca_nghiem)

# Keep only real roots using .is_real
nghiem_thuc = [n for n in tat_ca_nghiem if n.is_real]
print("Real roots:", nghiem_thuc)
\end{lstlisting}

\begin{lstlisting}[style=outstyle,caption={Kết quả}]
All roots: [2, -1 - sqrt(3)*I, -1 + sqrt(3)*I]
Real roots: [2]
\end{lstlisting}

\begin{infobox}
Dòng \texttt{[n for n in tat\_ca\_nghiem if n.is\_real]} là một
\textbf{list comprehension} (cú pháp rút gọn Python): duyệt qua từng
nghiệm \texttt{n}, chỉ giữ lại nghiệm mà \texttt{n.is\_real} trả về
\texttt{True}. Thuộc tính \texttt{.is\_real} là một trong nhiều
``thuộc tính suy luận'' (assumptions) SymPy tự động gắn cho mỗi biểu
thức: có thể trả về \texttt{True}, \texttt{False}, hoặc \texttt{None}
(nếu SymPy chưa đủ thông tin kết luận).
\end{infobox}

\section{Ví dụ phức tạp hơn -- ba số hạng có tích và tổng}

\begin{lstlisting}[caption={Bài toán: tổng và tích ba số hạng liên tiếp}]
u1, q = sp.symbols('u1 q', real=True)

# Given: u1 + u2 + u3 = 14  and  u1 * u2 * u3 = 64
pt1 = sp.Eq(u1 + u1*q + u1*q**2, 14)
pt2 = sp.Eq(u1 * (u1*q) * (u1*q**2), 64)

tap_nghiem = sp.nonlinsolve([pt1, pt2], [u1, q])
print(tap_nghiem)
\end{lstlisting}

\begin{lstlisting}[style=outstyle,caption={Kết quả (rút gọn)}]
{(2, 2), (8, 1/2)}
\end{lstlisting}

\section{Rút gọn nghiệm phức tạp với \texttt{simplify()}}

\begin{lstlisting}[caption={Rút gọn biểu thức nghiệm}]
bt_phuc_tap = (u1**2 * q**2 - u1**2 * q**0) / (u1 * q - u1)
print("Before:", bt_phuc_tap)
print("After: ", sp.simplify(bt_phuc_tap))
\end{lstlisting}

\begin{lstlisting}[style=outstyle,caption={Kết quả}]
Before: (u1**2*q**2 - u1**2)/(u1*q - u1)
After:  u1*(q + 1)
\end{lstlisting}

\begin{exbox}
\textbf{Bài tập 4.1.} Cho CSN thỏa $u_1 + u_3 = 20$ và $u_2 + u_4 = 40$.
Tìm $u_1, q$.\\[2pt]
\textbf{Bài tập 4.2.} Cho CSN có $u_2 = 6$ và $u_5 = 162$. Tìm $u_1$ và
$q$.\\[2pt]
\textbf{Bài tập 4.3.} Cho 3 số hạng liên tiếp có tổng bằng 21 và tích
bằng 216. Tìm ba số đó.\\[2pt]
\textbf{Bài tập 4.4.} Giải phương trình $q^4 - 5q^2 + 4 = 0$ và lọc
nghiệm thực dương.
\end{exbox}

\newpage

% =====================================================================
\chapter{Dạng toán 3 -- Cấp số nhân lùi vô hạn và giới hạn}

\section{Nhắc lại lý thuyết}

Khi $|q| < 1$, các số hạng của cấp số nhân càng lúc càng nhỏ dần về 0,
và tổng $S_n$ tiến dần đến một giá trị giới hạn hữu hạn $S$ khi
$n \to \infty$:
\[
S = u_1 + u_1 q + u_1 q^2 + \cdots = \frac{u_1}{1 - q}.
\]

\section{Tính giới hạn bằng \texttt{limit()}}

\begin{syntaxbox}{sp.limit(bieu\_thuc, bien, gia\_tri\_tien\_toi)}
Tính \textbf{giới hạn} của \texttt{bieu\_thuc} khi \texttt{bien} tiến
tới \texttt{gia\_tri\_tien\_toi}. Để biểu diễn ``tiến tới vô cùng'',
SymPy dùng ký hiệu \texttt{sp.oo} (đọc là \emph{infinity}, ký hiệu
$\infty$).
\end{syntaxbox}

\begin{lstlisting}[caption={Tính tổng CSN lùi vô hạn bằng limit}]
import sympy as sp
sp.init_printing()

u1, q, n = sp.symbols('u1 q n', real=True)
Sn_formula = u1 * (1 - q**n) / (1 - q)

# Limit of S_n as n -> infinity
S_vo_han = sp.limit(Sn_formula, n, sp.oo)
print(S_vo_han)
\end{lstlisting}

\begin{lstlisting}[style=outstyle,numbers=none]
zoo
\end{lstlisting}

\begin{warnbox}
Kết quả \texttt{zoo} xuất hiện vì SymPy \textbf{chưa biết} $q$ thỏa
điều kiện $|q| < 1$ -- về mặt toán học thuần túy, nếu $q$ có thể là
bất kỳ số thực nào (kể cả $q > 1$), giới hạn không tồn tại hữu hạn. Ta
cần \textbf{thay $q$ bằng giá trị cụ thể} trong khoảng $(-1, 1)$ trước
khi lấy giới hạn.
\end{warnbox}

\begin{lstlisting}[caption={Thay q cụ thể rồi lấy giới hạn}]
u1 = sp.symbols('u1', real=True)
q = sp.symbols('q', real=True)
n = sp.symbols('n')

Sn_formula = u1 * (1 - q**n) / (1 - q)

# Example: q = 1/2 (|q| < 1 guaranteed)
Sn_vi_du = Sn_formula.subs(q, sp.Rational(1, 2))
S_vo_han_vi_du = sp.limit(Sn_vi_du, n, sp.oo)
print("With q = 1/2, S =", S_vo_han_vi_du)
\end{lstlisting}

\begin{lstlisting}[style=outstyle,numbers=none]
With q = 1/2, S = 2*u1
\end{lstlisting}

\begin{infobox}
Kết quả $S = 2u_1$ khớp chính xác với công thức lý thuyết
$S = \dfrac{u_1}{1-q} = \dfrac{u_1}{1 - 1/2} = 2u_1$. Trong thực hành,
cách \textbf{đơn giản và an toàn nhất} là sau khi tự kiểm tra bằng tay
điều kiện $|q| < 1$ đúng theo đề bài, áp dụng thẳng công thức
$S = \dfrac{u_1}{1-q}$ (không cần gọi \texttt{limit()} nữa) -- hàm
\texttt{limit()} chủ yếu giúp \textbf{kiểm chứng} công thức xuất phát
từ định nghĩa toán học.
\end{infobox}

\section{Hàm tổng quát tính tổng CSN lùi vô hạn}

\begin{lstlisting}[caption={Hàm kiểm tra điều kiện và tính tổng}]
def tong_CSN_lui_vo_han(u1_val, q_val):
    """Compute S = u1/(1-q), requires |q| < 1."""
    if abs(q_val) >= 1:
        raise ValueError(
            f"No finite sum since |q|={abs(q_val)} >= 1"
        )
    u1s, qs = sp.symbols('u1 q')
    return (u1s / (1 - qs)).subs({u1s: u1_val, qs: q_val})

print(tong_CSN_lui_vo_han(4, sp.Rational(1, 3)))   # result: 6
\end{lstlisting}

\begin{lstlisting}[style=outstyle,numbers=none]
6
\end{lstlisting}

\section{Ví dụ tổng hợp -- bài toán chữ}

\begin{lstlisting}[caption={Tính tổng chuỗi hình học hội tụ}]
# S = 1 + 1/2 + 1/4 + 1/8 + ...
print(tong_CSN_lui_vo_han(1, sp.Rational(1, 2)))   # 2

# S = 3 - 1 + 1/3 - 1/9 + ...
print(tong_CSN_lui_vo_han(3, sp.Rational(-1, 3)))  # 9/4
\end{lstlisting}

\begin{exbox}
\textbf{Bài tập 5.1.} Tính $S = 5 + \tfrac{5}{4} + \tfrac{5}{16} +
\cdots$.\\[2pt]
\textbf{Bài tập 5.2.} Viết hàm kiểm tra một số $x$ có phải là công bội
hợp lệ cho CSN lùi vô hạn hay không.\\[2pt]
\textbf{Bài tập 5.3.} Cho $S = 8$ và $u_1 = 2$. Tìm $q$ (biết
$|q| < 1$).
\end{exbox}

\newpage

% =====================================================================
\chapter{Ứng dụng thực tế}

\section{Bài toán lãi suất kép}

Nếu gửi số tiền gốc $A$ với lãi suất $r$ mỗi kỳ, số dư sau $n$ kỳ tạo
thành một cấp số nhân với công bội $q = (1+r)$:
\[
T_n = A (1 + r)^n.
\]

\begin{lstlisting}[caption={Tính số dư sau n năm gửi tiết kiệm lãi kép}]
A, r, n = sp.symbols('A r n', positive=True)

# Balance after n periods at rate r per period
so_du = A * (1 + r)**n

# Deposit 100 million, 6%/year, after 10 years
so_du_thuc_te = so_du.subs({A: 100_000_000, r: sp.Rational(6, 100), n: 10})
print("Rounded (VND):", round(float(so_du_thuc_te.evalf()), 0))
\end{lstlisting}

\begin{lstlisting}[style=outstyle,numbers=none]
Rounded (VND): 179084769.0
\end{lstlisting}

\begin{infobox}
Cột mốc thời gian $n$ đóng vai trò giống hệt chỉ số số hạng trong cấp
số nhân đã học, còn $(1+r)$ đóng vai trò công bội $q$. Việc ``nhận ra''
một bài toán thực tế có bản chất là cấp số nhân (dân số tăng theo tỉ lệ
\% cố định, vi khuẩn nhân đôi sau mỗi giờ, khấu hao tài sản theo tỉ lệ
\% mỗi năm, \dots) chính là bước quan trọng nhất -- sau đó chỉ cần áp
dụng lại đúng \texttt{subs()}, \texttt{evalf()} đã học ở Dạng 1.
\end{infobox}

\section{So sánh lãi đơn và lãi kép}

\begin{lstlisting}[caption={So sánh lãi đơn (cấp số cộng) và lãi kép (cấp số nhân)}]
A, r, n = sp.symbols('A r n', positive=True)

lai_don = A * (1 + n * r)         # simple interest
lai_kep = A * (1 + r)**n          # compound interest

# Compare after 20 years with A = 100, r = 5%/year
A_val, r_val, n_val = 100, sp.Rational(5, 100), 20
print("Simple:  ", lai_don.subs({A: A_val, r: r_val, n: n_val}))
print("Compound:", lai_kep.subs({A: A_val, r: r_val, n: n_val}).evalf(6))
\end{lstlisting}

\begin{lstlisting}[style=outstyle,caption={Kết quả}]
Simple:   200
Compound: 265.330
\end{lstlisting}

\section{Bài toán tăng trưởng vi khuẩn}

\begin{lstlisting}[caption={So sánh tăng trưởng vi khuẩn với hai tốc độ}]
u1_val = 100    # initial 100 bacteria
q1 = sp.Rational(3, 2)   # +50% per hour
q2 = 2                   # doubles per hour

u1s, qs, ns = sp.symbols('u1 q n')
cong_thuc = u1s * qs**(ns - 1)

for gio in range(1, 6):
    sl1 = cong_thuc.subs({u1s: u1_val, qs: q1, ns: gio})
    sl2 = cong_thuc.subs({u1s: u1_val, qs: q2, ns: gio})
    print(f"Hour {gio}: +50% -> {sl1},  doubled -> {sl2}")
\end{lstlisting}

\begin{lstlisting}[style=outstyle,caption={Kết quả (rút gọn)}]
Hour 1: +50% -> 100,   doubled -> 100
Hour 2: +50% -> 150,   doubled -> 200
Hour 3: +50% -> 225,   doubled -> 400
Hour 4: +50% -> 675/2, doubled -> 800
Hour 5: +50% -> 5063/8, doubled -> 1600
\end{lstlisting}

\section{Bài toán dân số}

% ====== ĐÃ SỬA LỖI TẠI ĐÂY ======
% Dấu % trong caption phải được escape thành \%
\begin{lstlisting}[caption={Dự báo dân số với tỉ lệ tăng 1.5\%/năm}]
u1s, qs, ns = sp.symbols('u1 q n')
cong_thuc = u1s * qs**(ns - 1)

# Year 2025: 2 million, growth 1.5%/year, predict 2040
ket_qua = cong_thuc.subs({u1s: 2, qs: sp.Rational(15, 1000) + 1, ns: 15})
print("Population 2040 (million):", ket_qua.evalf(6))
\end{lstlisting}

\begin{lstlisting}[style=outstyle,numbers=none]
Population 2040 (million): 2.48454
\end{lstlisting}

\begin{exbox}
\textbf{Bài tập 6.1.} Gửi ngân hàng 500 triệu với lãi kép 7\%/năm. Sau
15 năm, số dư là bao nhiêu?\\[2pt]
\textbf{Bài tập 6.2.} Một tài sản trị giá 1 tỉ đồng khấu hao 12\%/năm.
Sau 8 năm giá trị còn bao nhiêu?\\[2pt]
\textbf{Bài tập 6.3.} Một quần thể vi khuẩn ban đầu có 200 con, sau 6
giờ có 12800 con. Tính công bội $q$ mỗi giờ.
\end{exbox}

\newpage

% =====================================================================
\chapter{Trực quan hóa với Matplotlib}

\section{Vẽ đồ thị so sánh các cấp số nhân}

\begin{lstlisting}[caption={Vẽ đồ thị so sánh các CSN với q khác nhau}]
import matplotlib.pyplot as plt
import sympy as sp

u1s, qs, ns = sp.symbols('u1 q n')
cong_thuc = u1s * qs**(ns - 1)

# Convert symbolic expression to fast numeric function
f = sp.lambdify((u1s, qs, ns), cong_thuc, 'numpy')

so_hang = list(range(1, 11))     # n = 1..10
danh_sach_q = [0.5, 1.2, 1.5, 2] # multiple ratios

plt.figure(figsize=(8, 5))
for q_val in danh_sach_q:
    gia_tri = [f(10, q_val, n_val) for n_val in so_hang]
    plt.plot(so_hang, gia_tri, marker='o', label=f"q = {q_val}")

plt.xlabel("Term index n")
plt.ylabel("Value u_n (with u1 = 10)")
plt.title("Growth comparison of geometric progressions")
plt.legend()
plt.grid(True, linestyle='--', alpha=0.5)
plt.tight_layout()
plt.savefig("growth_chart_example.png", dpi=150)
plt.show()
\end{lstlisting}

\begin{syntaxbox}{sp.lambdify(bien, bieu\_thuc, 'numpy')}
Chuyển một \textbf{biểu thức ký hiệu SymPy} thành một \textbf{hàm số
Python bình thường}, tính toán bằng số rất nhanh (dựa trên
\texttt{numpy}) -- hữu ích khi cần tính hàng trăm/nghìn giá trị để vẽ
đồ thị, vì gọi \texttt{.subs()} lặp lại nhiều lần sẽ chậm hơn đáng kể.
\end{syntaxbox}

\section{Đồ thị hàm mũ liên tục và dãy rời rạc}

\begin{lstlisting}[caption={So sánh dãy rời rạc và hàm mũ liên tục}]
import numpy as np
import matplotlib.pyplot as plt

q = 2
u1 = 1

# Discrete points
n_int = np.arange(1, 9)
u_int = [u1 * q**(k - 1) for k in n_int]

# Continuous exponential curve
x_cont = np.linspace(1, 8, 200)
y_cont = u1 * q**(x_cont - 1)

plt.figure(figsize=(8, 5))
plt.plot(x_cont, y_cont, '--', color='gray',
         label=r'Continuous $y = 2^{x-1}$')
plt.scatter(n_int, u_int, color='red', s=60, zorder=3,
            label=r'Discrete terms $u_n$')
plt.xlabel('n')
plt.ylabel('u_n')
plt.title('Geometric progression as a discrete exponential function')
plt.legend()
plt.grid(alpha=0.3)
plt.show()
\end{lstlisting}

\begin{tipbox}
Ba yếu tố \textbf{bắt buộc} của một đồ thị minh họa: \texttt{xlabel}
(nhãn trục hoành), \texttt{ylabel} (nhãn trục tung), \texttt{legend}
(chú thích phân biệt các đường), cùng \texttt{title} (tiêu đề) -- giúp
người xem hiểu ngay đồ thị biểu diễn điều gì mà không cần đọc lại code.
\end{tipbox}

\begin{infobox}
Quan sát đồ thị sẽ thấy: chỉ cần $q$ nhích từ $1.2$ lên $2$, độ chênh
lệch giữa các đường cong tăng \textbf{rất nhanh} sau vài số hạng --
đây chính là bản chất tăng trưởng theo cấp số nhân (exponential
growth), khác hẳn tăng trưởng đều (cấp số cộng).
\end{infobox}

\begin{exbox}
\textbf{Bài tập 7.1.} Vẽ đồ thị $u_n$ cho $u_1 = 100$, $q$ lần lượt
$0.9, 1.0, 1.1$ trong khoảng $n = 1..15$.\\[2pt]
\textbf{Bài tập 7.2.} Vẽ đồ thị so sánh số dư lãi kép và lãi đơn theo
$n = 1..30$ với $A = 100, r = 5\%$.
\end{exbox}

\newpage

% =====================================================================
\chapter{Đóng gói thành thư viện -- từ Notebook đến \texttt{.py}}

Sau khi các đoạn mã trong Notebook đã chạy đúng, bước cuối cùng là gom
các hàm lại thành một file Python độc lập, có thể \texttt{import} và tái
sử dụng ở bất kỳ dự án nào khác.

\section{Module \texttt{solver.py} hoàn chỉnh}

\begin{lstlisting}[caption={geometric\_progression\_solver.py -- đóng gói hàm}]
"""
Solver module for geometric progression using SymPy.
Provides: compute u_n, S_n, infinite sum, solve system for u1, q.
"""
import sympy as sp


def tinh_un(u1_val, q_val, n_val):
    """Return u_n = u1 * q^(n-1)."""
    u1s, qs, ns = sp.symbols('u1 q n')
    return (u1s * qs**(ns - 1)).subs({u1s: u1_val, qs: q_val, ns: n_val})


def tinh_Sn(u1_val, q_val, n_val):
    """Return S_n; handles the special case q = 1."""
    if q_val == 1:
        return n_val * u1_val
    u1s, qs, ns = sp.symbols('u1 q n')
    ct = u1s * (1 - qs**ns) / (1 - qs)
    return ct.subs({u1s: u1_val, qs: q_val, ns: n_val})


def tong_lui_vo_han(u1_val, q_val):
    """Return S = u1/(1-q), requires |q| < 1."""
    if abs(q_val) >= 1:
        raise ValueError("Not convergent: need |q| < 1")
    u1s, qs = sp.symbols('u1 q')
    return (u1s / (1 - qs)).subs({u1s: u1_val, qs: q_val})


def giai_he_tim_u1_q(pt1, pt2, chi_lay_nghiem_thuc=True):
    """Solve 2 equations (sp.Eq) for (u1, q). Returns list of solutions."""
    u1, q = sp.symbols('u1 q', real=chi_lay_nghiem_thuc)
    nghiem = sp.solve([pt1, pt2], [u1, q])
    return nghiem
\end{lstlisting}

\section{File kiểm thử \texttt{test\_solver.py}}

\begin{lstlisting}[caption={test\_solver.py -- kiểm thử đơn giản}]
from geometric_progression_solver import tinh_un, tinh_Sn, tong_lui_vo_han
import sympy as sp


def test_tinh_un():
    # u1 = 3, q = 2 -> u5 = 3 * 2^4 = 48
    assert tinh_un(3, 2, 5) == 48
    # u1 = 1, q = 1/2, n = 5 -> u5 = 1/16
    assert tinh_un(1, sp.Rational(1, 2), 5) == sp.Rational(1, 16)


def test_tinh_Sn_q_khac_1():
    # u1 = 3, q = 2 -> S6 = 189
    assert tinh_Sn(3, 2, 6) == 189


def test_tinh_Sn_q_bang_1():
    # q = 1 -> Sn = n * u1
    assert tinh_Sn(5, 1, 10) == 50


def test_tong_lui_vo_han():
    # u1 = 4, q = 1/2 -> S = 8
    assert tong_lui_vo_han(4, sp.Rational(1, 2)) == 8


if __name__ == "__main__":
    test_tinh_un()
    test_tinh_Sn_q_khac_1()
    test_tinh_Sn_q_bang_1()
    test_tong_lui_vo_han()
    print("All tests PASSED!")
\end{lstlisting}

\begin{tipbox}
Chạy file kiểm thử bằng \texttt{python test\_solver.py}, hoặc chuyên
nghiệp hơn bằng \texttt{pytest} (cài qua \texttt{pip install pytest}
rồi chạy \texttt{pytest test\_solver.py}).
\end{tipbox}

\newpage

% =====================================================================
\chapter*{Kết luận}
\addcontentsline{toc}{chapter}{Kết luận}

\section*{Tổng kết hành trình}
\addcontentsline{toc}{section}{Tổng kết hành trình}

Qua 9 chương của tài liệu, chúng ta đã đi từ những bước cài đặt môi
trường đơn giản nhất đến việc đóng gói hoàn chỉnh một module Python có
thể tái sử dụng cho việc giải toán cấp số nhân. Điểm then chốt cần ghi
nhớ:

\begin{enumerate}
  \item \textbf{Ký hiệu hóa} các biến bằng \texttt{symbols()} là bước
  đầu tiên để SymPy hiểu bài toán.
  \item \textbf{\texttt{subs()} và \texttt{evalf()}} là công cụ chính
  cho Dạng 1 (thay số tính toán).
  \item \textbf{\texttt{Eq()}, \texttt{solve()}, \texttt{nonlinsolve()}}
  giúp giải hệ phương trình tìm $u_1, q$ -- kèm việc \emph{lọc nghiệm
  thực}.
  \item \textbf{\texttt{limit()}} tự động hóa việc tính tổng lùi vô
  hạn mà không cần nhớ công thức.
  \item \textbf{Matplotlib} giúp trực quan hóa tốc độ tăng trưởng --
  kỹ năng quan trọng khi trình bày bài toán ứng dụng.
  \item \textbf{Đóng gói hàm} giúp chương trình tái sử dụng và kiểm
  thử được.
\end{enumerate}

\section*{Đánh giá hiệu quả}
\addcontentsline{toc}{section}{Đánh giá hiệu quả}

SymPy giúp giải quyết bài toán cấp số nhân theo đúng \textbf{ba lớp độ
khó tăng dần}: (1) thay số tính toán trực tiếp, (2) giải hệ phương trình
tìm ẩn, và (3) mô hình hóa các bài toán thực tế phức tạp hơn (lãi kép,
tăng trưởng, giới hạn). Việc giữ biểu thức ở dạng ký hiệu càng lâu càng
tốt (chỉ ``chốt'' về số thập phân ở bước cuối bằng \texttt{.evalf()})
giúp tránh sai số tích lũy và cho phép tái sử dụng công thức cho nhiều
bài tập khác nhau chỉ bằng một dòng \texttt{.subs()}.

\section*{Lưu ý và hạn chế}
\addcontentsline{toc}{section}{Lưu ý và hạn chế}

\begin{itemize}
  \item Hệ phương trình bậc cao có thể cho \textbf{nhiều nghiệm} (kể cả
  nghiệm phức) -- luôn cần đối chiếu với điều kiện thực tế của đề bài.
  \item Luôn xử lý riêng trường hợp $q = 1$ khi tính $S_n$ để tránh lỗi
  chia cho 0.
  \item Với cấp số nhân lùi vô hạn, phải tự đảm bảo điều kiện $|q| < 1$
  trước khi áp dụng công thức $S = \dfrac{u_1}{1-q}$.
  \item Nên dùng \texttt{sp.Rational()} thay vì số thập phân Python
  thường trong suốt quá trình tính toán, chỉ chuyển sang
  \texttt{.evalf()} ở bước trình bày kết quả cuối.
  \item \texttt{solve()} đôi khi bỏ sót nghiệm của hệ bậc cao -- cần
  đối chiếu với \texttt{nonlinsolve()}.
\end{itemize}

\section*{Hướng mở rộng}
\addcontentsline{toc}{section}{Hướng mở rộng}

\begin{itemize}
  \item Viết giao diện dòng lệnh (CLI) bằng \texttt{argparse} để chạy
  trực tiếp từ Terminal.
  \item Xây dựng web app bằng \texttt{Streamlit} để người dùng nhập đề
  bài và nhận lời giải.
  \item Bổ sung tính năng xuất lời giải ra LaTeX/PDF.
  \item Mở rộng sang cấp số cộng, dãy truy hồi, hoặc các dạng toán khác.
\end{itemize}

\section*{Tài liệu tham khảo}
\addcontentsline{toc}{section}{Tài liệu tham khảo}

\begin{enumerate}
  \item Sách giáo khoa Đại số và Giải tích -- Chuyên đề Dãy số, Cấp số
  cộng, Cấp số nhân (chương trình phổ thông Việt Nam).
  \item Tài liệu chính thức của SymPy: \url{https://docs.sympy.org}
  \item Tài liệu chính thức của Matplotlib: \url{https://matplotlib.org}
  \item Tài liệu Jupyter Notebook: \url{https://jupyter.org/documentation}
\end{enumerate}

% =====================================================================
\appendix

\chapter{Bảng tra cứu nhanh các lệnh SymPy}

\begin{longtable}{@{}p{4.2cm}p{4.2cm}p{6cm}@{}}
\toprule
\textbf{Lệnh} & \textbf{Cú pháp} & \textbf{Ý nghĩa}\\
\midrule
\endhead
Khai báo biến & \texttt{sp.symbols('x y')} & Tạo đối tượng ký hiệu\\
Biến có điều kiện & \texttt{sp.symbols('q', real=True)} & Ràng buộc miền giá trị\\
In đẹp & \texttt{sp.init\_printing()} & Output dạng LaTeX/Unicode\\
Thay giá trị & \texttt{expr.subs(\{x: 1\})} & Thay biến bằng giá trị\\
Số thực hóa & \texttt{expr.evalf(15)} & Chuyển sang số thập phân\\
Phân số chính xác & \texttt{sp.Rational(1, 2)} & Tạo phân số $\tfrac12$\\
Rút gọn & \texttt{sp.simplify(expr)} & Tối giản biểu thức\\
Khai triển & \texttt{sp.expand(expr)} & Mở ngoặc\\
Nhân tử & \texttt{sp.factor(expr)} & Phân tích thành nhân tử\\
Phương trình & \texttt{sp.Eq(A, B)} & Tạo phương trình $A = B$\\
Giải PT & \texttt{sp.solve(pt, x)} & Giải phương trình theo $x$\\
Giải hệ & \texttt{sp.solve([pt1, pt2], [x, y])} & Giải hệ phương trình\\
Hệ phi tuyến & \texttt{sp.nonlinsolve([...], [...])} & Giải hệ phi tuyến\\
Giới hạn & \texttt{sp.limit(expr, n, sp.oo)} & Giới hạn khi $n \to \infty$\\
Vô cùng & \texttt{sp.oo} & Ký hiệu $\infty$\\
Số ảo & \texttt{sp.I} & Đơn vị ảo $i$\\
Kiểm tra thực & \texttt{expr.is\_real} & True/False/None\\
Hàm số hóa & \texttt{sp.lambdify(x, expr)} & Chuyển sang hàm số nhanh\\
\bottomrule
\end{longtable}

\chapter{Một số lỗi thường gặp và cách khắc phục}

\begin{enumerate}
  \item \textbf{Lỗi \texttt{ZeroDivisionError} khi $q = 1$:}\\
  Nguyên nhân: công thức $S_n$ có mẫu $(1 - q)$.\\
  Cách xử lý: kiểm tra \texttt{if q == 1} trước khi áp dụng công thức.

  \item \textbf{Kết quả chứa \texttt{I} (số ảo) dù đề yêu cầu thực:}\\
  Nguyên nhân: SymPy giải trên trường số phức.\\
  Cách xử lý: khai báo \texttt{symbols('q', real=True)} hoặc lọc
  \texttt{not expr.has(sp.I)}.

  \item \textbf{Kết quả là biểu thức quá phức tạp:}\\
  Nguyên nhân: SymPy chưa rút gọn.\\
  Cách xử lý: gọi \texttt{sp.simplify(expr)}.

  \item \textbf{Sai số khi dùng \texttt{float}:}\\
  Nguyên nhân: \texttt{1/3} trong Python là \texttt{0.3333...} (số thực
  nhị phân).\\
  Cách xử lý: dùng \texttt{sp.Rational(1, 3)} thay vì \texttt{1/3}.

  \item \textbf{\texttt{solve()} không trả về nghiệm nào dù có nghiệm:}\\
  Nguyên nhân: hệ phi tuyến bậc cao.\\
  Cách xử lý: chuyển sang \texttt{nonlinsolve()}.

  \item \textbf{\texttt{limit()} trả về \texttt{zoo}:}\\
  Nguyên nhân: SymPy chưa biết điều kiện hội tụ $|q| < 1$.\\
  Cách xử lý: thay $q$ bằng giá trị cụ thể trong $(-1, 1)$ trước khi
  lấy giới hạn.

  \item \textbf{Nhầm lẫn giữa \texttt{=} và \texttt{==}:}\\
  Nguyên nhân: nhầm gán giá trị với so sánh.\\
  Cách xử lý: dùng \texttt{=} để gán biến; dùng \texttt{sp.Eq()} để
  biểu diễn phương trình.
\end{enumerate}

\chapter{Bài tập tổng hợp}

\section*{A. Dạng cơ bản (thay số)}

\begin{enumerate}
  \item Cho CSN có $u_1 = 4, q = 3$. Tìm $u_8$ và $S_8$.
  \item Cho $u_1 = 1, q = \tfrac{2}{3}$. Tính $S_{15}$ dạng phân số.
  \item Cho $u_1 = 1000, q = 1.02$. Sau 20 kỳ, $u_{20}$ là bao nhiêu?
\end{enumerate}

\section*{B. Dạng trung bình (giải hệ)}

\begin{enumerate}
  \setcounter{enumi}{3}
  \item Cho CSN thỏa $u_1 + u_3 = 30$ và $u_2 + u_4 = 60$. Tìm $u_1, q$.
  \item Cho CSN có $u_3 = 12$ và $u_6 = 96$. Tìm $u_1, q$.
  \item Ba số hạng liên tiếp của CSN có tổng bằng 14 và tổng nghịch đảo
  bằng $\tfrac{7}{8}$. Tìm ba số đó.
\end{enumerate}

\section*{C. Dạng nâng cao (ứng dụng)}

\begin{enumerate}
  \setcounter{enumi}{6}
  \item Gửi ngân hàng 200 triệu, lãi kép 8\%/năm. Sau 12 năm, số dư là
  bao nhiêu? Vẽ đồ thị tăng trưởng.
  \item Một quần thể thỏ ban đầu có 50 con, tăng trưởng 20\% mỗi năm.
  Dự báo số lượng sau 10 năm.
  \item Chứng minh bằng SymPy: nếu $u_k^2 = u_{k-1} \cdot u_{k+1}$ thì
  dãy $(u_n)$ là cấp số nhân.
  \item Viết hàm Python trả về $u_1, q$ khi cho trước $u_m, u_n, m, n$
  (bài toán tổng quát tìm CSN khi biết hai số hạng).
\end{enumerate}

\end{document}
